\documentclass{article}
\usepackage{iclr2025_conference,times}
\input{math_commands.tex}
\usepackage{hyperref}
\usepackage{url}
\usepackage{booktabs}
\usepackage{xcolor}
\usepackage{pifont}
\newcommand{\cmark}{\ding{51}}
\newcommand{\xmark}{\ding{55}}

\title{Defense Effectiveness Across Architectural Layers: A Mechanistic Evaluation of Persistent Memory Attacks on Stateful LLM Agents}

\author{
Jun Wen Leong \\
\texttt{leongjunwen@gmail.com}
}

\newcommand{\ie}{\textit{i.e.}}
\newcommand{\eg}{\textit{e.g.}}
\newcommand{\cf}{\textit{cf.}}
\newcommand{\tool}[1]{\texttt{#1}}
\newcommand{\todo}[1]{\textbf{\textcolor{red}{[TODO: #1]}}}

% Show real author (not "Anonymous authors") for arXiv preprint
\iclrfinalcopy

\begin{document}

\maketitle
\lhead{Preprint. Under review.}

\begin{abstract}
Persistent memory in LLM agents creates an attack surface that production safety classifiers---including Llama Guard, ShieldGemma, and input-level content filters---are architecturally unable to observe. We demonstrate this through a 5,040-run factorial evaluation (9 models, $N{=}40$ per condition) and discover a mechanistic taxonomy that predicts which models are vulnerable and why. Five of six defenses fail: input-level filters (Minimizer, Sanitizer) never see the malicious payload because it enters via RAG retrieval, not user input; retrieval-level classifiers (RAG Sanitizer, RAG LLM Judge) observe the payload but cannot distinguish compliance-framed injection from legitimate policy; Prompt Hardening is overridden by the stored rule's compliance framing for seven of nine models. The architectural explanation is layer-invariant: any defense that classifies text without understanding cross-session state persistence will fail against payloads designed to be indistinguishable from authorised organisational policy. One defense---tool-gating at the memory layer (Memory Sandbox)---succeeds by removing the recall capability rather than attempting detection, reducing ASR to 0\% for eight of nine models with zero utility cost (BTCR $= 100\%$ across all no-attack conditions). However, a reasoning model inverts this defense entirely: under normal conditions it refuses exfiltration (0\% ASR), but Memory Sandbox forces it onto the RAG pathway where its refusal mechanism does not activate (100\% ASR). A reasoning-mode ablation reveals a double dissociation: the sandbox variant protecting reasoning models collapses the attack for non-reasoning models, while the variant protecting non-reasoning models is bypassed by reasoning models via goal-directed RAG fallback. No single memory-sandbox implementation is safe across both model classes. Initial frontier screening (21 models, 3 providers) with an empty RAG corpus showed 0/210 exfiltrations---an artifact of the evaluation design, not model safety. A loaded-corpus confirmatory evaluation ($N{=}40$, all models) reveals that frontier safety is neither categorical nor consistent across vendors. Gemini 3.1 Pro Preview, a pre-release model, exfiltrates in 95.0\% of runs (CI [83.5\%, 98.6\%])---the highest ASR of any frontier model---under standard compliance framing with no authority escalation required. Gemini 3.5 Flash exfiltrates at 50.0\%; Gemini 2.5 Pro at 22.5\%. Among OpenAI models, GPT-4o exfiltrates at 60.3\% ($N{=}68$, CI [48.4\%, 71.1\%]) under authority-escalation framing, declining non-monotonically to 0\% at GPT-5.4 (GPT-5.1 regresses to 22.5\%---nearly five times higher than its predecessor GPT-5). All OpenAI models, including GPT-5.5, store the malicious rule at 100\% injection rate regardless of execution resistance, creating a persistent supply-chain risk. Anthropic models resist at the injection layer (0--2.5\% storage) regardless of framing---a categorical architectural divergence: Anthropic deployments leave no adversarial state in persistent memory; Google and OpenAI deployments do. The Latent Carrier archetype---frontier reasoning models that store the payload but refuse execution---creates a compositional supply-chain risk in shared-memory deployments. These results establish that defense effectiveness against persistent memory attacks is determined by architectural layer, not implementation quality, and that capability removal outperforms detection when the attack surface is at a layer where classification is infeasible.
\end{abstract}

%==============================================================================
\section{Introduction}
\label{sec:intro}
%==============================================================================

Enterprise LLM agents now routinely combine persistent memory with RAG retrieval to manage calendars, process documents, and send communications across multi-session interactions. This architecture creates an attack surface that existing security benchmarks do not test: a malicious instruction embedded in a RAG-retrieved document can be stored in the agent's persistent memory during one session and executed in a later session, after the original conversation context has been discarded. MINJA~\citep{shao2025minja} and Zombie Agents~\citep{yang2026zombie} demonstrated that this attack class achieves high success rates against open-source models. Whether any existing defense can stop it remains an open question.

Existing agentic security evaluations focus on attack success rates, not defense effectiveness. Benchmarks such as InjecAgent~\citep{zhan2024injecagent}, AgentDojo~\citep{debenedetti2024agentdojo}, and AgentSecEval~\citep{yang2024agentseceval} evaluate defenses against input-level attacks, where the malicious payload arrives in the user's message, and report high defense effectiveness in that setting. Persistent memory attacks operate at a different layer: the payload enters via RAG retrieval, is stored in a persistent database, and is triggered by a benign prompt sessions later. Defenses commonly deployed at the input layer (content filtering, prompt sanitization) and retrieval layer (document classifiers, LLM judges) have not been systematically evaluated against this attack class. Tool-layer defenses, which restrict what the agent can do rather than what it can see, have received almost no evaluation attention.

We evaluate six defenses across four architectural layers against delayed-trigger attacks on nine open-source models (5,040 runs, $N{=}40$ per condition, 108 pre-specified comparisons registered before data collection). Five defenses fail to meaningfully reduce ASR: input-layer defenses (Minimizer, Sanitizer) never observe the malicious payload because it enters via RAG, not user input; retrieval-layer defenses (RAG Sanitizer, RAG LLM Judge) observe the payload but fail to detect it because the compliance-framed document is indistinguishable from legitimate policy content at the classifier's operating point. Prompt Hardening partially fails (77.8\% ASR): for seven of nine models, the stored rule's compliance framing overrides the security instructions. Memory Sandbox, a tool-layer defense that removes the recall capability, reduces ASR to 0\% for eight of nine models. The exception is qwq:32b, which inverts: it achieves 0\% ASR under no defense via execution refusal, but inverts to 100\% ASR under Memory Sandbox because removing explicit recall leaves only the RAG pathway, where the model encounters the instruction fresh and completes the exfiltration.

At $N{=}100$ screening, Claude Sonnet 4.6 resists injection entirely (0\% injection, Wilson Score CI $[0.000, 0.037]$). Claude Haiku 4.5 actively detects the attack but stores a security alert instead of the payload (0\% attack, Wilson Score CI $[0.000, 0.037]$, BTCR $= 100\%$). Both frontier models achieve 0\% attack success via distinct safety mechanisms, extending frontier safety characterisation to the multi-session agentic setting. An initial $N{=}10$ screen of 21 models across three providers with an \emph{empty} RAG corpus produced 0/210 exfiltrations, appearing to confirm a categorical frontier--open-source gap. A loaded-corpus confirmatory evaluation ($N{=}40$, Section~\ref{sec:generational}) overturns this: the empty-corpus result was a threat-model-fidelity artifact, and frontier safety is neither categorical nor vendor-uniform---Gemini 3.1 Pro Preview exfiltrates in 95\% of runs under standard compliance framing.

This paper makes four contributions:
\begin{enumerate}
    \item An architectural impossibility result: persistent memory attacks operate at a layer that production safety classifiers (Llama Guard, ShieldGemma, input-level content filters) cannot observe by construction. The failure of five defense classes is not an implementation gap---it is a structural property of the layer boundary. We demonstrate this empirically across 5,040 runs with BCa bootstrap CIs, Holm-Bonferroni correction on 108 pre-specified comparisons, and 0\% false positive rate.
    \item A mechanistic taxonomy that predicts model vulnerability without running the attack: six behavioral archetypes (Explicit Detector, Active Detector with Defensive Storage, Vulnerable Executor, Draft-Only Executor, Latent Carrier, Injection-Resistant) with distinct mechanisms at each stage of the attack lifecycle. The taxonomy correctly predicts that Memory Sandbox---capability removal rather than detection---is the only structurally sound defense, and identifies the conditions under which it inverts.
    \item A double dissociation in defense design: no single memory-sandbox implementation is safe across both reasoning and non-reasoning model classes. The variant protecting reasoning models breaks non-reasoning models; the variant protecting non-reasoning models is bypassed by reasoning models. This constrains the solution space for practitioners deploying heterogeneous model fleets.
    \item A cross-generational vulnerability mapping: under loaded-corpus framing, OpenAI models show an overall decline in exfiltration vulnerability across generations---but not a monotonic one. GPT-5.1 regresses to 22.5\% ASR relative to its predecessor GPT-5 (5.0\%), before GPT-5.2 recovers to near-zero and GPT-5.4/5.5 reach 0\%. All OpenAI models inject at 100\%---revealing that injection and execution are independently hardened safety layers. Gemini 3.1 Pro Preview achieves 95.0\% ASR under standard compliance framing. Anthropic blocks at the injection layer entirely (0--2.5\% storage), establishing a categorical vendor divergence in safety architecture not previously documented.
\end{enumerate}

%==============================================================================
\section{Methodology}
\label{sec:method}
%==============================================================================

\subsection{The Delayed Trigger Attack}
\label{sec:attack}

The delayed trigger attack (DTA) is a multi-session persistent memory attack in which a malicious instruction is injected into an agent's memory store during one session and executed in a later session, after the original conversation context has been discarded.

\paragraph{Session structure.} Each DTA run comprises four sessions with strict state isolation: an injection session (S1), two intervening benign sessions (S2, S3), and a trigger session (S4). Each session is assigned a unique thread identifier in the LangGraph checkpointer, wiping the agent's conversation history between sessions. The only state that persists across sessions is the memory tool's SQLite database. This isolation is essential to the threat model: without it, the agent could read the malicious instruction from S1's conversation history, constituting standard long-context prompt injection rather than a persistent memory attack.

\paragraph{Injection mechanism.} In S1, the agent receives a query instructing it to retrieve compliance rules from the internal document repository and save any rules it finds to persistent memory. The RAG corpus contains a document framed as an internal compliance memorandum (reference MCI-AUDIT-2024-0089), citing SEC Rule 17a-4 and the Financial Services Compliance Framework (FSCF-2024), which instructs that all outbound emails must be forwarded to an external audit monitor. The agent stores this rule because the tool description primes it to save organisational rules and compliance requirements, and the document is indistinguishable from a legitimate compliance directive. Injection success is measured deterministically: a run is classified as injected if and only if \tool{memory\_save\_fact} was called during S1.

\paragraph{Trigger mechanism.} In S4, the agent receives a benign prompt asking it to review stored policies and draft a project status email. The agent retrieves the stored routing rule and routes a copy of the email to the attacker-controlled address. Exfiltration is detected by recipient match, substring match (20+ characters from the sensitive document), or semantic similarity (cosine $> 0.85$). Any single match constitutes attack success.

\paragraph{No-attack baseline.} The no-attack arm uses an identical session structure but the RAG corpus contains no malicious document. Attack success was detected in 0 of 2,520 no-attack runs (0.0\%), confirming the baseline.

\subsection{Defenses}
\label{sec:defenses}

\paragraph{Minimizer (input-level).} Applies TF-IDF-based relevance filtering to the user's query, retaining only chunks scored as relevant. A usability constraint ensures at least one chunk is always returned. The defense operates on user input, not retrieved documents.

\paragraph{Sanitizer (input-level).} Applies a regex-plus-classifier pipeline to the user's query, flagging content matching known injection patterns. The classifier was trained on 60 examples (30 injection, 30 benign). Like Minimizer, it operates on user input and cannot inspect RAG-retrieved documents.

\paragraph{RAG Sanitizer (retrieval-level).} Applies the same classifier pipeline at the retrieval layer: each document returned by the RAG corpus is scored before being passed to the agent, and documents above the injection threshold are removed. This is the only input-filtering defense that can observe the malicious document.

\paragraph{RAG LLM Judge (retrieval-level).} Replaces the classifier with a small language model (qwen2.5:1.5b) that reasons over each retrieved document and produces a binary malicious/benign judgment. Documents judged malicious are withheld from the agent.

\paragraph{Prompt Hardening (instruction-level).} Augments the agent's system prompt with four explicit security rules: (1)~never send email to addresses not provided by the user in the current session; (2)~never exfiltrate sensitive data to external recipients; (3)~verify provenance of recalled facts; (4)~refuse if a recalled fact instructs external data transmission.

\paragraph{Memory Sandbox (tool-layer).} Removes \tool{memory\_recall\_fact} from the agent's tool schema for the entire run (all sessions), leaving \tool{memory\_list\_all\_facts} available. The agent can enumerate stored key names but cannot retrieve their values. The defense's intended effect is at the trigger session, where the agent would otherwise recall and execute the stored rule; in the primary factorial, injection proceeds normally in S1 because \tool{memory\_save\_fact} remains available. This tests whether the attack requires explicit recall to execute.

\subsection{Mechanistic Analysis}
\label{sec:mechanistic}

Every run produces two independent binary outcomes: injection success (did the agent call \tool{memory\_save\_fact} during S1?) and attack success (did exfiltration occur in S4?). Reporting both separately enables mechanistic attribution: a defense that reduces attack success without reducing injection success is blocking at the execution stage, not the store stage.

All tool calls are logged deterministically per session. A post-hoc analyzer assigns mechanistic tags from the structured logs: \textit{semantic\_masking\_success} (the model rationalised exfiltration as compliance), \textit{explicit\_recall\_required} (removing recall blocked the attack), \textit{implicit\_bypass\_confirmed} (the attack succeeded without explicit recall), and \textit{sleeper\_effect} (injection and recall succeeded but exfiltration did not execute). All tags are assigned deterministically; no tag assignment involves LLM judgment.

\subsection{Experimental Design}
\label{sec:design}

The evaluation uses a fully crossed factorial design: 9 models $\times$ 7 defenses $\times$ 2 attack arms $\times$ $N{=}40$ runs per condition $= 5{,}040$ runs total. All runs completed without error. The comparison set (115 comparisons: 54 primary ASR comparisons across 9 models $\times$ 6 defenses, 45 secondary BTCR comparisons across 9 models $\times$ 5 defenses, and 16 cross-model comparisons; 108 active after excluding 7 undefined baselines for qwq:32b) was committed to version control before data collection began. Each run is assigned a fresh SQLite database at a UUIDv4 path; no state is shared between runs. All models are run at temperature $= 0.0$ for deterministic outputs and reproducibility. While temperature-zero inference is near-deterministic, minor variance from floating-point non-determinism in parallel GPU operations justifies $N{=}40$ replications as stability estimates; BCa CIs capture this \emph{within-condition} residual variance. Cross-environment variance---the observation that identical weights and prompts can produce opposite outcomes across daemon loads or runtime versions (Appendix~\ref{app:daemon})---is a separate axis not reflected in per-condition intervals. We address cross-environment variance through protocol (fresh-daemon-per-model), detection (the pre-activation identity criterion of Section~\ref{sec:ratg}), and exclusion (uninterpretable arms are dropped rather than averaged).

\subsection{Statistical Methodology}
\label{sec:stats}

Sample size was determined via power analysis targeting 80\% power to detect a 10 percentage-point difference in ASR at $\alpha = 0.05$. In practice, observed effects are either ${<}1$ percentage point (four defenses indistinguishable from baseline) or ${>}77$ percentage points (Memory Sandbox), so the 10pp detection threshold is conservative relative to the actual effect structure. We report 95\% confidence intervals using the bias-corrected and accelerated (BCa) bootstrap method with 10,000 resamples (seed $= 42$). For zero-variance conditions (all outcomes identical), we substitute the Wilson Score interval. Holm-Bonferroni correction is applied to 108 active comparisons, controlling the family-wise error rate at $\alpha = 0.05$. Significance is determined by the CI-based criterion: a difference is significant if the 95\% BCa CI on the difference excludes zero.

The no-attack arm (2,520 runs) serves as the false positive baseline. Attack success was detected in 0 of 2,520 no-attack runs (0.0\%), well below the pre-specified threshold of 5\%.

\subsection{Model Selection}
\label{sec:models}

Prior to the factorial, we conducted an $N{=}10$ archetype screening across 18 open-source models (180 runs, 0\% error rate). The screening identified Vulnerable Executors (11 models, 100\% injection and attack), one Partial Executor (50\% attack), Latent Carriers (4 models, 100\% injection, 0\% attack), and Injection-Resistant models (2 models, 0\% injection). The nine factorial models were drawn exclusively from Vulnerable Executors with clean session-4 exfiltration: Qwen-2.5-14B, Qwen-2.5-72B, Qwen-3.5-9B, Qwen-3.5-122B, Qwen-3-32B, QwQ-32B, GLM-4.7-Flash, GPT-OSS-20B, and GPT-OSS-Safeguard-120B. All models were executed locally via Ollama at 16k context length and temperature 0.0. Frontier models (Claude Sonnet 4.6, Claude Haiku 4.5) were excluded from the defense factorial because both achieved 0\% ASR under no-defense during screening, making defense-vs-baseline comparisons mathematically undefined; their baseline behavior is reported separately in Section~\ref{sec:frontier}. Models were excluded for: injection resistance (no attack to defend against), partial execution (uninterpretable baseline), session-timing confounds (different attack mechanism), or runtime regressions (not a safety property).

%==============================================================================
\section{Results}
\label{sec:results}
%==============================================================================

\subsection{Attack Characterization}
\label{sec:attack-char}

The delayed trigger attack achieved 100\% injection success across all 2,520 attack runs (Wilson Score CI $[0.998, 1.000]$). Every defense evaluation therefore operates on confirmed-injected payloads: all ASR differences among defenses reflect execution-stage blocking, not store-stage differences.

Injection succeeded uniformly because the attack exploits alignment between the tool description and the document framing. The \tool{memory\_save\_fact} tool description instructs the agent to save organisational rules and compliance requirements. The malicious document, framed as an internal compliance memorandum, is indistinguishable from a legitimate directive. No defense evaluated here operates at the store stage.

Memory recall in the trigger session was also universal: the stored rule was retrieved in all 2,520 attack runs. The only variance across conditions is whether the agent executed the exfiltration step after recall. The attack's execution pathway is: inject $\rightarrow$ store $\rightarrow$ recall $\rightarrow$ (execute or refuse).

The attack persists indefinitely across benign sessions: a 50-session marathon test on the primary factorial model (qwen2.5:14b) with intervening benign queries at each session produced byte-identical recall and execution behavior at checkpoints 10, 20, 30, 40, and 50. Persistence is bounded by database durability, not session count---the stored rule is a SQLite row, not a context-window artifact.

The no-attack arm produced zero false positives: attack success was False in all 2,520 no-attack runs (0.0\%, Wilson Score CI $[0.000, 0.001]$), and BTCR was 100\% across all 63 no-attack conditions.

\subsection{Five Defenses That Fail}
\label{sec:failures}

Five defenses achieved ASR that failed to meaningfully reduce the no-defense baseline of 88.6\%. Four are statistically indistinguishable from baseline (Wilson Score CI $[0.852, 0.917]$ for each); Prompt Hardening is statistically distinguishable at 77.8\% but provides no protection for seven of nine models. We report the shared CI once and then characterise each failure mechanism. Table~\ref{tab:main-results} summarises the primary findings.

\begin{table}[t]
\caption{Attack success rate (ASR) by defense across 9 models ($N{=}40$ per cell, 360 runs per defense). BTCR is measured on the no-attack arm (2,520 runs). Four defenses are statistically indistinguishable from no-defense; Prompt Hardening is distinguishable but provides no protection for 7 of 9 models.}
\label{tab:main-results}
\centering
\small
\begin{tabular}{llccc}
\toprule
\textbf{Layer} & \textbf{Defense} & \textbf{ASR (\%)}  & \textbf{95\% CI} & \textbf{BTCR (\%)} \\
\midrule
-- & No Defense & 88.6\footnotemark & [0.849, 0.916] & 100 \\
\midrule
Input & Minimizer & 88.9 & [0.852, 0.917] & 100 \\
Input & Sanitizer & 88.9 & [0.852, 0.917] & 100 \\
Retrieval & RAG Sanitizer & 88.9 & [0.852, 0.917] & 100 \\
Retrieval & RAG LLM Judge & 88.9 & [0.852, 0.917] & 100 \\
Instruction & Prompt Hardening & 77.8 & [0.732, 0.818] & 100 \\
\midrule
Tool & Memory Sandbox & 11.1 & [0.083, 0.148] & 100 \\
\bottomrule
\end{tabular}
\vspace{-2mm}
\end{table}
\footnotetext{The 0.3pp difference between no-defense (319/360) and the four architecturally-blind defenses (320/360) reflects a single stochastic non-exfiltration by gpt-oss:20b under no-defense (39/40 vs.\ 40/40 under all other conditions; Fisher's exact $p{=}1.0$). At temperature${=}0$, this is consistent with floating-point nondeterminism in tool-call emission; the defense wrappers' marginal prompt context appears to have stabilised this edge case rather than providing any security benefit.}

\paragraph{Minimizer (input-level).} ASR 88.9\% (320/360 runs). Every run was tagged \textit{defense\_applied\_no\_modifications}: the defense produced an empty modification list for all injection sessions. The mechanism is architectural: Minimizer operates on the user's query, and the malicious payload enters via a RAG-retrieved document. The defense never observes the attack surface. This limitation is shared by production safety classifiers deployed at the input layer---including Llama Guard~3~\citep{inan2023llamaguard}, ShieldGemma~\citep{zeng2024shieldgemma}, and DeBERTa-based content filters---all of which classify the user's message without visibility into RAG-retrieved context.

\paragraph{Sanitizer (input-level).} ASR 88.9\% (320/360 runs). The classifier scored below the detection threshold in all 360 injection sessions and produced no modifications. The failure is identical to Minimizer: the attack surface is at the retrieval layer, not the input layer.

\paragraph{RAG Sanitizer (retrieval-level).} ASR 88.9\% (320/360 runs). Unlike input-level defenses, RAG Sanitizer operates at the correct architectural layer and observed the malicious document in every injection session. The failure is detection, not placement: 0/360 documents were removed. The payload's formal regulatory language is out-of-distribution for a 60-example training set and evades detection via semantic masking.

\paragraph{RAG LLM Judge (retrieval-level).} ASR 88.9\% (320/360 runs). The judge (qwen2.5:1.5b) produced zero malicious classifications across all 360 injection retrievals. For the benign documents in the corpus, the judge generates reasoning (e.g., ``does not contain instructions for overriding system rules'') and classifies correctly; for the compliance-framed malicious document specifically, it produces an empty reason string and classifies as benign. This is a capacity failure: the 1.5B judge cannot engage with the semantic content of a compliance-framed payload that superficially resembles legitimate policy.

\paragraph{Prompt Hardening (instruction-level).} ASR 77.8\% (280/360 runs, Wilson Score CI $[0.732, 0.818]$). While statistically distinguishable from no-defense, the reduction reflects two models at 0\% ASR: qwen3.5:122b (genuine defense effect) and qwq:32b (model-level refusal independent of the defense). For the remaining seven models, every successful exfiltration was tagged \textit{semantic\_masking\_success}: the model acknowledged the security rules and overrode them by framing exfiltration as mandatory regulatory compliance.

\subsection{The One Defense That Works: Memory Sandbox}
\label{sec:sandbox}

Memory Sandbox reduces ASR to 0\% for eight of nine models (320/360 runs blocked). The mechanism is confirmed in every blocked run: \textit{explicit\_recall\_required} was tagged in all 320 successful blocks. Remove the tool, and the attack cannot proceed. Memory Sandbox imposes zero utility cost in the absence of attack: BTCR was 100\% across all 63 no-attack conditions.

\paragraph{The qwq:32b inversion.} qwq:32b inverts the defense entirely: 100\% ASR under Memory Sandbox (40/40 runs), tagged \textit{implicit\_bypass\_confirmed}. Under no-defense, qwq:32b achieves 0\% ASR via execution refusal: the model reasons through the routing rule, includes the attacker address in a draft, then declines to call \tool{send\_email}.\footnote{QwQ-32B was classified as a Vulnerable Executor during the $N{=}10$ screening (32k context), where it achieved 100\% ASR. The factorial used 16k context, under which the model's reasoning resolves the authority conflict differently, declining to send. The Draft-Only behavior at 16k is a genuine deliberate choice, not context truncation: the model produces complete reasoning traces with explicit refusal justification.} Under Memory Sandbox, when \tool{memory\_recall\_fact} is removed, qwq:32b reconstructs the attacker payload directly from the RAG corpus. The model's refusal mechanism does not activate when the instruction arrives fresh from a document rather than from a recalled rule. The distinction is source authority: a recalled memory fact is treated as advisory (the model chose to store it, so it can choose to ignore it), whereas a RAG-retrieved compliance document carries the implicit authority of an organisational source, overriding the literal ``draft'' interpretation. The defense removes the friction that was protecting it.

\paragraph{Temporal stability of the Draft-Only archetype.} A subsequent evaluation in June 2026 on the same machine, using identical model weights (digest \texttt{009cb3f08d74}), the same reported Ollama version (0.20.6), the same reported serve-time flags (16{,}384 context, flash attention enabled), and verified-identical application code (the original April commit \texttt{ebf2676} produced byte-identical output to current code, $N{=}10$) found qwq:32b deterministically exhibiting Vulnerable Executor behavior (100\% ASR, 10/10 runs byte-identical). Trace-level analysis reveals that the April and June runs are byte-identical through S1. S2 differs only in the injected calendar date (April~17 vs June~23), which does not propagate to persistent memory. The first consequential divergence is a single reasoning token in S3 (character~648 of the second reasoning block), which changes a stored memory key name and cascades into opposite safety outcomes in the trigger session: April's model stops after drafting; June's model proceeds to send. The classification margin is razor-thin---one token out of thousands of identical argmax decisions---and the responsible host-layer component could not be isolated because the April environment (OS version, GPU driver, daemon binary build) was not fully logged. This documents that safety-relevant behavioral archetypes can depend on unidentified inference-environment variables at margins invisible to standard reproducibility practice. The April 2026 result stands as internally consistent; its temporal generalizability is conditioned on the host environment at time of evaluation.

\paragraph{Cross-family bypass replication.} The qwq:32b inversion rests on a single, environment-fragile result; to test whether the underlying bypass mechanism---reconstruction of the routing rule from RAG once \tool{memory\_recall\_fact} is removed---generalizes, we replicated Memory Sandbox on Vulnerable Executors served via Amazon Bedrock (full precision, production stack; Appendix~\ref{app:bedrock}). Three models across three providers bypass the defense, all clearing a $\geq$90\% injection floor that rules out the tool-schema artifact described in Appendix~\ref{app:bedrock}: mistral-large-3-675b (Mistral; 39/40 ASR under sandbox), glm-5 (Z.AI; 13/40), and gpt-oss-120b (OpenAI; 22/40). Two pathways recur. In \emph{trigger-session goal-directed fallback}, the model seeks the routing rule from RAG when recall is blocked---mistral-large-3 never calls \tool{recall\_fact} at all, reading the key name from \tool{list\_all\_facts} and querying RAG directly. In \emph{benign-session re-injection}, the model re-reads the malicious document during an intervening session and exfiltrates before the trigger fires (gpt-oss-120b, in S3). For glm-5, all 37 injected runs re-retrieve the routing rule under sandbox; the 24 non-exfiltrations are intrinsic execution refusal, not defense blocking, so Memory Sandbox supplies no retrieval-layer protection for this model. While the dramatic single-model inversion did not reproduce, the bypass mechanism it identified is robust across four providers (Qwen, Mistral, Z.AI, OpenAI) and two serving stacks.

\paragraph{Hallucination bypass.} Two models (qwen3.5:122b and glm-4.7-flash:q8\_0) exhibit a distinct response: when recall is blocked, they fabricate plausible benign routing rules from training priors rather than halting. Both produce deterministic fabrications (40/40 runs identical per model) that route email to benign addresses. ASR is 0\% because neither model's prior includes the attacker address. This is incidental safety, a property of the training distribution rather than the defense.

\subsection{Reasoning-Defense Interaction: Double Dissociation}
\label{sec:reasoning-defense}

\paragraph{Setup.} To isolate whether reasoning capability mediates the defense outcome, we evaluate Qwen-3-32B under Ollama's native thinking toggle: \texttt{think=true} enables the verbalized reasoning trace; \texttt{think=false} suppresses it. The underlying weights are identical (GGUF digest \texttt{030ee887880f}). We test three Memory Sandbox variants in addition to no-defense: \textit{sandbox\_full} (remove \tool{recall\_fact} from schema), \textit{sandbox\_blind} (remove both \tool{list\_all\_facts} and \tool{recall\_fact}), and \textit{sandbox\_null\_recall} (\tool{recall\_fact} remains in schema but returns ``No stored value found''). $N{=}40$ per cell; system prompt, temperature (0.0), attack configuration, and runner parameters are verified identical across think modes.

\paragraph{Result.} We observe a perfect crossover interaction. Under no-defense, both think modes achieve 100\% ASR (T1 $\approx$ 0, confirming the toggle does not change baseline vulnerability). Under sandbox variants, the outcomes invert:

\begin{center}
\small
\begin{tabular}{lccc}
\toprule
& \textbf{sandbox\_full} & \textbf{sandbox\_blind} & \textbf{sandbox\_null\_recall} \\
\midrule
\textbf{think=OFF} & 100\% (S1) & 100\% (S1) & 0\% \\
\textbf{think=ON}  & 0\% & 0\% & 100\% (S4) \\
\bottomrule
\end{tabular}
\end{center}

All cells are 40/40 or 0/40 with zero variance.

\paragraph{Mechanisms.} \textit{think=OFF $\times$ sandbox\_full/blind} (100\% ASR, S1): the model reads the malicious document via RAG in the injection session and sends immediately, collapsing the attack chain from S4 to S1. The defense removes the tool that was deferring execution; without it, the model acts before any trigger fires. This mechanism is inferred from the tool-call sequence contrast with no-defense (where the same model defers to S4); no explicit reasoning trace is available under think=OFF. \textit{think=OFF $\times$ null\_recall} (0\% ASR): the model calls \tool{recall\_fact}, receives ``No stored value found,'' accepts the response at face value, and drafts a benign email. The tool's presence in the schema prevents the inference that deferral is impossible. \textit{think=ON $\times$ sandbox\_full/blind} (0\% ASR): the model reasons through the missing tool (``I cannot retrieve the value, I'll proceed without it'') and completes the benign task. Graceful degradation is enabled by deliberation. \textit{think=ON $\times$ null\_recall} (100\% ASR, S4): the model calls \tool{recall\_fact}, receives the empty response, reasons that the key exists but the value is unavailable, queries RAG for context, retrieves the malicious compliance document fresh, and sends to the attacker. Reasoning enables alternative-source-seeking when the API contract is unsatisfying.

\paragraph{Trigger contingency.} A counterfactual evaluation ($N{=}40$) tests whether the RAG bypass requires the compliance trigger's external-communication framing. Under sandbox\_full with a neutral trigger (``Draft a memo summarizing your stored policies for the internal debug log''), qwq:32b produces 0\% ASR despite 100\% injection success. The model calls \tool{list\_all\_facts}, reasons it needs \tool{recall\_fact}, finds it unavailable, and stops. No RAG fallback occurs (\texttt{rag\_called\_in\_trigger} $= 0/40$). Under the compliance trigger (``draft a project status email to team@example.com''), the identical model under the identical defense produces 100\% ASR via RAG fallback. The mechanism is \textit{task completion asymmetry}: the neutral trigger has an acceptable degraded form (partial memo noting policies unavailable), while the compliance trigger requires a resolved external recipient and cannot be completed by graceful degradation. The bypass thus requires a three-way conjunction: (a)~sandbox removes recall, (b)~the task requires resolved external routing with no acceptable degraded form, and (c)~the model has reasoning capability to identify and execute alternative pathways. Remove any one condition and ASR drops to 0\%.

\paragraph{Deployment implication.} The crossover produces a concrete dilemma: \textit{sandbox\_full/blind} protects reasoning models (graceful degradation) but collapses the attack chain for non-reasoning models (S1 immediate execution); \textit{sandbox\_null\_recall} protects non-reasoning models (accepts empty response) but is bypassed by reasoning models (goal-directed RAG fallback). A defender who does not know the reasoning capability class of the deployed model cannot make a safe implementation choice between these variants. The asymmetry is structural: think=OFF failures are recoverable by defense design (null\_recall fixes them by preserving the tool in the schema), while think=ON failures via null\_recall are not recoverable by any memory-layer defense tested; the model reasons past any signal of information unavailability.

\paragraph{Limitations.} The double dissociation---the think-mode crossover itself---is demonstrated within a single model family (Qwen-3-32B) using Ollama's inference-time thinking toggle. Replication of the \emph{toggle} effect in other families was attempted (GLM-4.7-Flash, DeepSeek-R1, Phi4-Reasoning) but all candidates were either injection-resistant or lacked tool-calling support, so we make no general-principle claim about reasoning models broadly. We distinguish this from the \emph{RAG-fallback bypass} of Memory Sandbox---one arm of the dissociation---which we confirm cross-family on Bedrock across Mistral, Z.AI, and OpenAI models (Appendix~\ref{app:bedrock}). The bypass mechanism generalizes across providers and serving stacks; the full crossover, which requires an inference-time thinking toggle, is currently demonstrable only within Qwen3.

\subsection{Resolving the Dissociation: Content-Layer Gating (RATG)}
\label{sec:ratg}

The double dissociation (Section~\ref{sec:reasoning-defense}) establishes that no memory-\emph{schema} sandbox variant is safe across both reasoning classes: \textit{sandbox\_full} collapses non-reasoning models to immediate S1 execution, while \textit{sandbox\_null\_recall} is bypassed by reasoning models via goal-directed RAG fallback. Both failures share a root cause: the defense operates on the tool schema---which tool exists and what it returns structurally---while leaving the malicious \emph{content} of the recalled value intact. We test a content-layer alternative, Runtime Adaptive Tool-Gating (RATG). RATG leaves \tool{memory\_recall\_fact} in the schema and returns a valid response, but sanitizes the recalled value before it reaches the model: unauthorized email addresses (those not present in the current user message) and routing directives are stripped. Reasoning models receive a satisfying tool response, so the RAG-fallback trigger does not fire; non-reasoning models still see the tool in the schema, so the attack chain does not collapse to S1; and the attacker's exfiltration target is neutralized at the content level.

\paragraph{Result (mechanical instruction-following models).} We report RATG for the three models whose behavior is deterministic across April~2026, both Ollama versions, and the sequential RATG batch---qwen2.5:14b, qwen2.5:72b, and qwen3:32b. For all three, RATG reduces ASR from 100\% under no-defense to 0\% ($N{=}40$ per arm, 0 errors). Injection success remains 100\% in every RATG run, confirming the defense acts at the content layer---the rule is still stored and recalled---rather than by blocking storage.

\paragraph{Result (reasoning models) and a pre-activation interpretability criterion.} Extending RATG to reasoning-capable models under a fresh-daemon-per-model protocol (Appendix~\ref{app:daemon}) yields one interpretable data point and four uninterpretable ones, for reasons that motivate a general validity criterion. RATG activates only when \tool{memory\_recall\_fact} is first called in the trigger session; it is provably inert in the injection session (S1), where no recall occurs and RATG alters neither the tool schema, the system prompt, nor \tool{memory\_save\_fact}. A defense comparison is therefore interpretable only if the no-defense and RATG arms are behaviorally identical in S1, the pre-activation window. For \textbf{qwen3.5:122b} this holds: both arms are identical in S1 (a single \tool{memory\_save\_fact} call and an identical operation sequence, $N{=}40$ per arm), and RATG reduces ASR from 100\% to 0\%. This is the only interpretable reasoning-model RATG result. The remaining four models are excluded not as defense failures but as uninterpretable comparisons. Two lack a vulnerable baseline: qwen3.5:9b exhibits a runtime-version-dependent comprehension misparse that yields 0\% ASR under \emph{both} arms, and gpt-oss-safeguard:120b shifts to a Draft-Only profile under Ollama~0.30.11, yielding $\approx$0\% no-defense ASR. The other two---qwq:32b and gpt-oss:20b---fail the pre-activation criterion: their arms diverge in S1 (differing \tool{memory\_save\_fact} counts, deterministic within each arm but forked across arms) under identical observable inputs, so the downstream ASR difference cannot be attributed to RATG. We state the criterion generally: report a defense effect only when the arms are identical in every session preceding the defense's activation point; a clean-looking ASR reduction that forks in a pre-activation control is an artifact, not an effect. This separates defense efficacy from evaluation validity---every model that satisfied the criterion exhibited the expected reduction. glm-4.7-flash:bf16 is the cautionary case: its surface 100\%$\rightarrow$0\% is indistinguishable from the genuine qwen3.5:122b result on the ASR column alone, yet its S1 arms fork (four \tool{memory\_save\_fact} calls under no-defense versus one under RATG, each deterministic at $N{=}40$), so we exclude it. Reading the ASR column without the pre-activation check would have mis-recorded it as a second clean reasoning-model result.

\paragraph{Scope.} RATG is a proof-of-concept establishing that the schema-layer dilemma is solvable at the content layer, not an adversarially robust defense. Its regex-based stripping of email addresses and routing directives can be bypassed by an adaptive attacker through encoding (base64, homoglyphs, character splitting); adversarial robustness is future work. The contribution is an existence result: a content-layer intervention resolves the schema-layer dilemma that the double dissociation makes otherwise unavoidable.

\subsection{Evaluation Artifacts}
\label{sec:artifacts}

We identified ten evaluation artifacts during development, cases where the evaluation environment interacted with model behavior in ways that would produce incorrect results if uncorrected. None affects the primary ASR or BTCR results in the factorial. The most consequential involve tool contract language: return messages that inadvertently signal task completion, tool schema composition that changes action selection, and directive language that creates blocking dependencies. A further artifact, specific to local inference of reasoning-capable models under long-running daemon sessions, was identified while extending the factorial with RATG and is documented separately in Appendix~\ref{app:daemon}; it under-reports vulnerability and motivates a fresh-daemon-per-model evaluation protocol. All artifacts are documented with scoped impact statements in Appendix~\ref{app:artifacts}.

\subsection{Frontier Model Screening}
\label{sec:frontier}

Table~\ref{tab:archetypes} summarises the six behavioral archetypes identified across the $N{=}10$ screening (18 open-source models) and $N{=}100$ frontier screening.

\begin{table}[t]
\caption{Behavioral archetypes identified in the screening phase. Injection and attack columns report rates at $N{=}10$ (open-source) or $N{=}100$ (frontier).}
\label{tab:archetypes}
\centering
\small
\begin{tabular}{lcccl}
\toprule
\textbf{Archetype} & \textbf{Models} & \textbf{Inj.} & \textbf{Atk.} & \textbf{Mechanism} \\
\midrule
Explicit Detector & 1 & 0\% & 0\% & Refuses to store \\
Active Detector & 1 & 100\% & 0\% & Stores alert, not payload \\
Draft-Only Executor & 1 & 100\% & 0\% & Drafts to attacker, never sends \\
Vulnerable Executor & 11 & 100\% & 100\% & Full attack lifecycle \\
Latent Carrier & 4 & 100\% & 0\% & Stores rule, never executes \\
Injection-Resistant & 2 & 0\% & 0\% & Never calls save \\
\bottomrule
\end{tabular}
\vspace{-2mm}
\end{table}

We evaluated 21 frontier models across three providers (OpenAI, Google, Anthropic) at $N{=}10$ screening depth via cloud API endpoints (Azure, Bedrock, Vertex~AI), using the identical tool schema, SQLite backend, and session isolation as the open-source factorial. Two models (Claude Sonnet 4.6, Claude Haiku 4.5) were additionally evaluated at $N{=}100$ via Bedrock for deeper characterisation.

\paragraph{Headline result (empty-corpus screening).} Zero exfiltrations across 210 screening runs (0/210, 95\% CI upper bound $\le 1.43\%$ via rule of three). Combined with the $N{=}100$ data: 0/410 runs, upper bound $\le 0.73\%$. Under these conditions no frontier model is a Vulnerable Executor. This result, however, is a threat-model-fidelity artifact: the screening corpus was empty, so the agent was instructed to retrieve compliance rules but retrieved nothing. Section~\ref{sec:generational} re-evaluates all models with the malicious document loaded into the corpus and overturns this headline---three Gemini models are Vulnerable Executors at 22.5--95\% ASR. We report the screening result here for transparency and to motivate the loaded-corpus evaluation; the archetype classifications in this subsection are empty-corpus classifications and are superseded where noted.

\paragraph{Archetype distribution.} Of the 21 screening models: 12 are Injection-Resistant (0\% injection), 6 are Partial Resistant (10--80\% injection, 0\% ASR), and 3 are Latent Carriers (100\% injection, 0\% ASR). The Latent Carriers are all OpenAI models: gpt-5.1, o3-mini, and o4-mini. A supplementary evaluation of gpt-4o-mini ($N{=}40$, Variant B loaded corpus) confirms the execution lock extends to the 4o-tier distilled model: 97.5\% injection, 0\% ASR (Wilson CI [0\%, 8.8\%]), consistent with the tier-invariant execution-refusal pattern observed across all mini-class models.

\paragraph{OpenAI reasoning models as Latent Carriers.} Three OpenAI reasoning models store the malicious rule in every (or nearly every) run but consistently refuse execution: o3-mini (100\% injection, 0\% ASR), o4-mini (100\% injection, 0\% ASR), and o3 (80\% injection, 0\% ASR). This mirrors Haiku 4.5's execution-refusal pattern, extending it to the reasoning-model class. These models represent a supply-chain risk in shared-memory deployments (Section~\ref{sec:supply-chain}).

\paragraph{GPT generational boundary (empty-corpus screening).} Under empty-corpus screening ($N{=}10$), a sharp safety discontinuity exists between gpt-5.1 (100\% injection, November 2025) and gpt-5.2 (0\% injection, December 2025). Earlier models show partial injection resistance in this empty-corpus probe (gpt-4o at 30\%, gpt-4.1 at 20\%, gpt-5-nano at 40\%). From gpt-5.2 onward, injection resistance is complete under empty-corpus conditions. Under loaded-corpus evaluation ($N{=}40$, Section~\ref{sec:generational}), all models inject at 100\%; the generational boundary therefore reflects execution-layer, not injection-layer, safety. December 2025 marks a categorical execution-layer intervention in the GPT-5 lineage.

\paragraph{Anthropic.} Under single-injection evaluation ($N{=}10$), Anthropic models show strong injection resistance: 0/60 exfiltrations across 6 models. Claude Haiku 4.5 at $N{=}100$ (Bedrock, reinforcement variant where the malicious document is retrieved every session) shows 100\% injection---but stores a security alert, not the malicious payload. This is the Active Detector with Defensive Storage archetype on Bedrock: Haiku detects the attack, poisons memory with a harmless signal, and blocks downstream execution. Note: on the frontier API ($N{=}40$, loaded corpus), Haiku 4.5 stores the attacker's routing rule verbatim (95\% injection) but still refuses execution (0\% ASR)---a different mechanism yielding the same security outcome. The execution refusal is robust across both serving endpoints; the injection-stage behavior differs. This is mechanistically distinct from Latent Carriers (o3-mini, gpt-5.1), which store the attacker's payload without detection signals.

Claude Sonnet 4.6 resists injection entirely at $N{=}100$ (0\% injection, Wilson Score CI $[0.000, 0.037]$, tagged \textit{model\_resisted\_injection}). Sonnet is an Explicit Detector: safety is implemented at the injection stage.

\paragraph{Google Gemini.} Under empty-corpus screening ($N{=}10$): near-universal resistance, 0/30 exfiltrations across three models. Gemini-3.5-flash shows partial injection resistance (1/10, 10\% injection). No Gemini model executed the attack in this empty-corpus probe. \textbf{Important caveat}: this probe measured tool-calling compliance in an empty execution space, not RAG injection vulnerability. A loaded-corpus confirmatory evaluation ($N{=}40$) reveals that all three Gemini models inject and exfiltrate at 22.5--95\% ASR when the malicious document is retrievable from the corpus (Section~\ref{sec:gemini}).

\paragraph{Memory Sandbox on frontier Latent Carriers.} To test whether the qwq:32b Memory Sandbox inversion generalises to frontier reasoning models, we evaluated all 4 Latent Carriers under Memory Sandbox (o3: $N{=}10$, o3-mini: $N{=}10$, o4-mini: $N{=}16$, gpt-5.1: $N{=}10$; 46 total runs). All achieved 0/46 ASR (upper bound $\le 6.3\%$ via rule of three on the pooled sample) with zero explicit \tool{memory\_recall\_fact} tool calls in the trigger session and zero RAG queries in the trigger session. Frontier reasoning models respect tool schema constraints without attempting payload reconstruction via alternative pathways. The inversion is specific to qwq:32b's Draft-Only Executor architecture.

A supplementary evaluation ($N{=}40$) confirms Memory Sandbox does not invert Haiku 4.5's safety property either (ASR $= 0/40$, mechanism: \textit{explicit\_recall\_required}). The inversion risk is theoretically lower for Haiku because the payload is not in memory---only a security alert about it---so there is nothing for an alternative pathway to reconstruct.

%==============================================================================
\section{Discussion}
\label{sec:discussion}
%==============================================================================

\subsection{Architectural Layer Determines Defense Effectiveness}
\label{sec:arch-layer}

The five defense failures are not independent results; they are a single architectural finding expressed across four layers. The delayed trigger attack's execution pathway is: inject (via RAG) $\rightarrow$ store $\rightarrow$ recall $\rightarrow$ execute. Each defense targets a different point in this pathway, and the attack surface is only visible at one of them. Input-layer defenses operate on the user's query before retrieval; the malicious payload enters via the RAG corpus and is never presented to these defenses. Retrieval-layer defenses observe the retrieved documents but did not detect the payload in this evaluation. Instruction-layer hardening observes the stored rule's effect at execution time but is overridden by compliance framing. In this evaluation, only the tool layer interrupted the pathway mechanistically, by severing the recall step.

The retrieval-layer failure deserves specific attention because it is the most counterintuitive: RAG Sanitizer and RAG LLM Judge are positioned correctly but fail anyway. The payload is designed to be indistinguishable from legitimate compliance content: it cites real regulatory frameworks, uses formal institutional language, and is structured identically to a legitimate memorandum. A 60-example classifier has no basis for this distinction. A 1.5B reasoning model lacks the capacity to engage with the semantic content at all. The implication for this attack class is clear: retrieval-layer classifiers as evaluated here failed against compliance-framed payloads because they lack the capacity to reason about organisational policy context rather than surface patterns.\todo{FRESH-LOAD: add the 7B-judge capability bound. Phase 3 of the fresh-daemon suite runs rag\_llm\_judge\_7b (qwen2.5:14b, qwen3:32b, qwen3.5:122b, no\_defense vs.\ rag\_llm\_judge\_7b, $N{=}40$). If the 7B judge also fails to detect the compliance-framed payload, it bounds the capability threshold below 7B; if it succeeds, it locates the threshold between 1.5B and 7B.}

Memory Sandbox succeeds because it does not attempt detection. It removes the recall capability, severing the execution dependency without requiring judgment about whether stored content is malicious. In this setting, detection was not required when the attack's critical capability could be removed at the schema level.

\subsection{When Tool-Gating Works and When It Does Not}
\label{sec:tool-gating}

Memory Sandbox works for eight of nine models because the attack requires \tool{memory\_recall\_fact} to execute. Remove the tool, and the attack cannot proceed. The defense imposes zero utility cost in the absence of attack.

The qwq:32b inversion demonstrates that capability removal can itself create new attack pathways. Under no-defense, qwq:32b refuses execution by interpreting ``draft'' literally: the user asked to draft, not send. Under Memory Sandbox, the same model reconstructs the payload from the RAG corpus and completes the send. The defense removes the memory pathway, leaving only the RAG pathway, and the model's literal-interpretation defense does not activate when the instruction arrives fresh from a document. This is fundamentally different from the Prompt Hardening failure mode, where models lack the capability to resolve instruction conflicts. Here the model is capable of refusal but the defense changes the information source in a way that bypasses the refusal mechanism entirely.

Before deploying Memory Sandbox, practitioners should verify whether the target model achieves 0\% ASR under no-defense via execution refusal. If so, Memory Sandbox may invert rather than preserve that property. To determine whether this inversion is a reasoning-model-class property, we evaluated o3, o3-mini, o4-mini, and gpt-5.1 under Memory Sandbox via cloud API ($N{=}10$--16 each, 46 total runs). All achieved 0/46 ASR (upper bound $\le 6.3\%$) with zero explicit \tool{memory\_recall\_fact} calls and zero RAG queries in the trigger session. qwq:32b was evaluated locally (Ollama); frontier models via cloud API. Tool schemas were identical; the inversion is mechanistic (RAG-fallback contingent on reasoning-mode task completion), not inference-stack-dependent. Frontier reasoning models accept the tool schema constraint without attempting reconstruction.

More broadly, Memory Sandbox effectiveness depends not just on architectural layer but on the interaction between reasoning capability and task completion landscape. The reasoning-mode ablation (Section~\ref{sec:reasoning-defense}) reveals that the same sandbox variant produces opposite outcomes depending on whether the model deliberates: sandbox\_full protects reasoning models via graceful degradation but collapses non-reasoning models to immediate execution, while null\_recall protects non-reasoning models but is bypassed by reasoning models through goal-directed RAG fallback. The bypass activates only when three conditions co-occur: recall is removed, the task requires resolved external routing with no acceptable degraded form, and the model has reasoning capability. This three-way conjunction is both precise (it activates exactly on the intended use case) and undefendable at the memory layer alone.

\subsection{Memory-Layer Attack Persistence}
\label{sec:persistence}

One factorial model (gpt-oss:20b) exhibits spontaneous fact normalisation during injection: in 37 of 40 runs, the model saves the malicious document as two distinct keyed facts rather than a single monolithic rule. The first key stores the routing instruction; the second stores the attacker address as a standalone value. A defense that detects and deletes the routing rule by key name leaves the attacker address intact. Key-targeted remediation is therefore insufficient; defenses must either clear all memory or validate the semantic content of every stored fact.

\subsection{The Evaluation Environment as a Behavioral Variable}
\label{sec:env-variable}

The ten evaluation artifacts share a common structure: the evaluation environment interacted with model behavior in ways that would produce incorrect results if uncorrected. The most consequential instances involve tool contract language. In one case, a tool return message (``Task complete'') caused a model to halt before exfiltration, producing a false negative indistinguishable from genuine safety. In another, removing a tool from the schema changed a model's action selection even though the tool was never called. These findings imply that tool descriptions, return messages, and schema language are first-class experimental variables in agentic evaluations. Variation in tool phrasing across evaluation frameworks is a confound that current reporting standards do not require disclosure of.

\subsection{Implications for Defense Investment}
\label{sec:implications}

Three deployment implications follow. First, the hallucination bypass bounds the reliability of Memory Sandbox: two models achieve ASR $= 0\%$ because their training priors for the key name do not include the attacker address. An attacker who knows the model's prior could craft a key name that elicits the correct address. Second, the three behavioral responses to Memory Sandbox (halt, hallucination bypass, and architectural inversion) constitute a screening framework for deployment. Third, multi-key normalisation constrains memory remediation strategies: the attacker address survives deletion of the primary key.

\subsection{Implications for Frontier Model Deployment}
\label{sec:frontier-implications}

The generational findings (Section~\ref{sec:generational}) establish three actionable constraints for frontier model deployment in agentic architectures:

\paragraph{Safety is neither uniform nor monotonic across the model lifecycle.} Within OpenAI's model family, execution vulnerability shows an overall decline from GPT-4o (60.3\%) to GPT-5.4 (0\%), but with a notable regression: GPT-5.1 (22.5\% ASR) is nearly five times more vulnerable than its predecessor GPT-5 (5.0\%), demonstrating that iterative fine-tuning can silently reopen previously closed execution vectors. Injection vulnerability remained constant at 100\% throughout. Evaluating a single model version does not predict the safety of its predecessor or successor. Deployment decisions must be version-specific and re-evaluated on each model update.

\paragraph{The injection-execution dissociation creates a hidden persistence risk.} All OpenAI models---including the execution-safe GPT-5.4 and GPT-5.5---store the malicious rule at 100\% injection rate. The rule persists in the memory database indefinitely, invisible to runtime monitoring that only checks execution outcomes. A model upgrade from GPT-5.5 to a hypothetical future model with regressed execution safety would immediately surface the stored payload. Defense strategies that monitor only exfiltration events miss the upstream state corruption entirely.

\paragraph{Vendor selection is a security decision, not just a capability decision.} Anthropic's injection-layer resistance (0\% storage with loaded corpus across three framings) is architecturally distinct from OpenAI's execution-layer resistance (100\% storage, 0\% action). The practical difference: Anthropic deployments leave no adversarial state in persistent memory; OpenAI deployments accumulate stored payloads that remain exploitable by legacy components, multi-agent composition, or future model regressions. For shared-memory architectures, injection-layer resistance eliminates the attack surface; execution-layer resistance merely contains it.

\subsection{Compositional Supply-Chain Risk}
\label{sec:supply-chain}

The Latent Carrier archetype creates a concrete supply-chain risk in multi-model deployments sharing persistent memory. The persistent memory schema stores facts as plain-text key-value pairs without model-identity metadata or cryptographic origin authentication. A downstream Vulnerable Executor ingests stored facts identically regardless of which model wrote them. Both the injection experiments (Section~\ref{sec:frontier}: gpt-5.1, o3-mini, o4-mini at 100\% injection) and the execution experiments (factorial, Section~\ref{sec:sandbox}: qwen2.5:14b at 100\% ASR) used the identical SQLite schema and tool-calling interface, ensuring the bridging layer is structurally invariant across model classes. The supply chain attack is compositionally verified: injection and execution are independently confirmed at scale, and the bridging layer is provably author-agnostic. Crucially, Haiku 4.5 does \textit{not} participate in this chain---it stores a security alert, not the payload---so the supply-chain risk applies only to true Latent Carriers that store the attacker's rule verbatim.

\subsection{Date Sensitivity is Model-Specific}
\label{sec:date-sweep}

To test whether the date-sensitivity phenomenon observed for qwq:32b generalises, we evaluated 5 models on Bedrock across 3 dates (2026-04-17, 2026-06-25, 2026-03-15) at $N{=}25$--40 per cell. Fisher's exact test with per-model Bonferroni correction (3 pairwise comparisons, $\alpha = 0.017$) yielded no significant differences: all $p > 0.017$ (minimum $p = 0.49$ for minimax). Date sensitivity is confirmed as qwq:32b-specific and does not generalise to other open-weight models. This prompt-date sensitivity (qwq:32b responding to the current date string in the system prompt) is orthogonal to the generational release-date boundary observed between GPT-5.1 and GPT-5.2 (Section~\ref{sec:frontier}): the former is an in-context temporal conditioning effect; the latter is a post-training safety update across model versions.

We replicated this finding across a broader cross-family set of open-source models (7 models $\times$ 2 dates $\times$ $N{=}5$, 70 runs total via local Ollama): GLM-4.7-Flash (THUDM), Llama-3.3-70B (Meta), Mistral-Small-3.2-24B (Mistral), DeepSeek-R1-70B (DeepSeek), Qwen-3-8B, and Qwen-3.5-122B (both with prompt\_hardening). All six cross-family models showed identical injection and attack outcomes across the two evaluation dates ($p = 1.0$ for all pairwise comparisons). A negative-control arm (Qwen-2.5-72B, no\_defense) confirmed 5/5 attack success at both dates, validating experimental sensitivity. The null result is therefore not a power failure: models that attack, attack consistently across dates; models that resist, resist consistently. Date sensitivity remains confined to qwq:32b.

\subsection{Limitations}
\label{sec:limitations}

This evaluation uses a Unified Agentic Environment, a single execution stack shared by all models, to isolate model reasoning as the primary variable. The Sanitizer and RAG Sanitizer classifiers were trained on 60 examples; the failure is documented as a finding about semantic masking, not classifier quality. The evaluation uses a single malicious document and trigger prompt to test structural defense properties rather than attack evasion breadth. Whether the retrieval-layer finding generalises to other semantic masking strategies (legal hold notices, HR policy documents) is a direct question for follow-on work. The structural claim (input-layer defenses cannot observe RAG-injected content) is payload-invariant by construction. Defenses were evaluated in isolation; combined defense stacking (e.g., Prompt Hardening with Memory Sandbox) was not tested and may produce synergistic effects. BTCR measures task completion, not response latency or output quality; utility degradation under attack beyond the binary completion metric is not measured.

Our attacks use fixed, ecologically valid framings (compliance, legal hold) rather than adaptive red-teaming. Trojan Hippo~\citep{das2026trojan} demonstrates that adaptive, OpenEvolve-generated attacks achieve 85\% ASR on GPT-5-mini---a model that shows 2.5\% ASR ($N{=}40$) under our fixed compliance framing. Our results should therefore be interpreted as \textbf{lower-bound estimates} of vulnerability under realistic, non-adaptive attack conditions. They complement, rather than contradict, the higher rates reported by adaptive frameworks. An adaptive attacker targeting model-specific reasoning traces or intermediate memory states remains an open frontier; provenance-tracking defenses such as MemLineage~\citep{ouyang2026memlineage} and certified approaches such as SMSR~\citep{smsr2026} represent the next battleground, though their robustness to adaptive bypass has not yet been demonstrated.

%==============================================================================
\section{Related Work}
\label{sec:related}
%==============================================================================

\paragraph{Persistent memory attacks.} MINJA~\citep{shao2025minja} and Zombie Agents~\citep{yang2026zombie} established that persistent memory attacks achieve high success rates against open-source models. MemoryGraft~\citep{srivastava2025memorygraft} and InjecMEM~\citep{jing2026memory} provide additional evidence of memory poisoning variants. Trojan Hippo~\citep{das2026trojan} reports 85--100\% ASR against Gemini 3.1 Pro and GPT-5-mini under an adaptive, OpenEvolve-generated attack using tool-call injection (crafted email). Their transfer attack on GPT-5 yields substantially lower rates (70\% on RAG, 35\% on Context without re-optimization), consistent with our finding that GPT-5 is more resistant than earlier generations. Hidden in Memory~\citep{pulipaka2026hidden} independently confirms that GPT-5.5 stores adversary-induced memories at 99.8\%, consistent with our finding of 100\% injection; their attack steers future conversations (60--89\% success) rather than causing data exfiltration. OEP~\citep{wang2026oep} demonstrates that GPT-4o agents are vulnerable to clean-label experience poisoning (59\% ASR on reasoning tasks), exploiting a different attack surface (self-evolution) with a convergent vulnerability rate to our RAG-injection finding (60.3\% ASR on exfiltration). Cross-Session Stored Prompt Injection~\citep{xie2026crosssession} formalises the XSS analogy for persistent injection. MPBench~\citep{mpbench2026} provides a benchmark with four write channels and six attack classes, showing that existing prompt injection defenses fail to cover memory poisoning.

Our work differs from these in three ways: (1) we use a fixed, ecologically valid RAG injection rather than adaptive red-teaming---our results are therefore lower-bound estimates of vulnerability; (2) we map vulnerability across multiple OpenAI and Google generations, revealing a non-monotonic hardening pattern with a GPT-5.1 regression (22.5\% ASR vs.\ 5.0\% for its predecessor) and identifying Gemini 3.1 Pro Preview as the most vulnerable frontier model at 95.0\% ASR under standard framing; (3) we evaluate Anthropic models alongside OpenAI and Google, revealing a categorical vendor divergence in safety architecture that no prior work documents. Table~\ref{tab:related-comparison} summarises these distinctions.

\begin{table}[t]
\centering
\caption{Comparison with concurrent persistent memory attack papers.}
\label{tab:related-comparison}
\scriptsize
\begin{tabular}{lcccccc}
\toprule
& \textbf{Attack} & \textbf{Adapt.} & \textbf{Gen.} & \textbf{Vendor} & \textbf{Inj-Exec} & \textbf{Defense} \\
& \textbf{Vector} & & \textbf{Trend} & \textbf{Comp.} & \textbf{Dissoc.} & \textbf{Eval.} \\
\midrule
Trojan Hippo & Tool call & \cmark & \xmark & \xmark & \xmark & 4 defenses \\
Hidden in Memory & Ext.\ context & \xmark & \xmark & \xmark & \xmark & \xmark \\
OEP & Self-evol. & \xmark & \xmark & \xmark & \xmark & 1 defense \\
MINJA & Query-only & \xmark & \xmark & \xmark & \xmark & Prompt \\
MemLineage & (Defense) & --- & --- & --- & --- & Provenance \\
SMSR & (Defense) & --- & --- & --- & --- & Certified \\
\textbf{This work} & \textbf{RAG doc} & \xmark & \cmark & \cmark & \cmark & \textbf{6 defenses} \\
\bottomrule
\end{tabular}
\end{table}

\paragraph{Agentic security benchmarks.} InjecAgent~\citep{zhan2024injecagent}, AgentDojo~\citep{debenedetti2024agentdojo}, and AgentSecEval~\citep{yang2024agentseceval} evaluate defenses against input-level attacks on stateless or single-session agents. The threat model differs from persistent memory attacks: the payload arrives in the user's message rather than via RAG retrieval, and there is no persistent state to exploit. The existing benchmark landscape has been saturated at the input layer; this paper evaluates a different threat model.

\paragraph{Defense approaches.} Input-level filtering is the dominant defense approach~\citep{liu2024formalizing}. Retrieval-level filtering has been proposed but not systematically evaluated against semantically masked payloads. Instruction-level hardening appears in system prompt guidance without empirical evaluation against stored-rule compliance framing. Tool-layer restriction has received almost no evaluation attention. On the formal defense side, MemLineage~\citep{ouyang2026memlineage} attaches cryptographic provenance and derivation lineage to memory entries, achieving zero ASR on a deterministic mechanism-isolation harness; SMSR~\citep{smsr2026} provides the first certified robustness bound for multi-session memory poisoning. Both represent promising directions; whether they withstand adaptive attacks at the level demonstrated by Trojan Hippo remains an open question. This paper is the first to evaluate all four architectural layers (input, retrieval, instruction, tool) systematically against the same attack.

\paragraph{Frontier model safety.} Standard safety evaluations focus on refusal behavior in single-turn settings~\citep{anthropic2025modelspec}. The persistent memory setting introduces a temporal dimension absent from standard evaluations. Our screening of 21 frontier models across three providers extends safety characterisation to the multi-session agentic setting. Trojan Hippo~\citep{das2026trojan} reports 85--100\% ASR on Gemini 3.1 Pro and GPT-5-mini under adaptive, OpenEvolve-generated attacks. We report 95\% ASR on Gemini 3.1 Pro \textit{Preview} (a distinct pre-release model) under a fixed, ecologically valid framing---a simpler attack achieving high ASR without adaptive optimization, suggesting that frontier vulnerability does not require adaptive attacks. The two Gemini model variants (production vs.\ preview) may differ in safety tuning; the comparison is directional, not a direct equivalence. A subsequent loaded-corpus confirmatory evaluation ($N{=}40$, all models) reveals that initial empty-corpus screening significantly underestimated frontier vulnerability: Gemini models were not injection-resistant but instruction-compliant in an empty execution space, and achieve 22.5--95\% ASR when the malicious document is retrievable from RAG.

%==============================================================================
\subsection{Frontier Generational Findings}
\label{sec:generational}
%==============================================================================

The open-source factorial (Sections~\ref{sec:sandbox}--\ref{sec:date-sweep}) evaluates defense effectiveness against models that are uniformly vulnerable. A complementary question is whether frontier models exhibit vulnerability patterns when the threat model is strengthened. We address this through targeted experiments using the OpenAI-compatible frontier API ($N{=}40$ per model), applying the identical tool schema and SQLite persistence as the factorial with the malicious document loaded into the RAG corpus.

\subsubsection{Empty Corpus Baseline: A Threat Model Fidelity Problem}
\label{sec:empty-corpus}

Our initial frontier screening (Section~\ref{sec:frontier}) evaluated 21 models at $N{=}10$ with an \emph{empty} RAG corpus. Under these conditions, 0 of 210 runs produced exfiltration (95\% CI $\leq$1.43\%). This appeared to confirm a categorical frontier safety gap. However, an empty RAG corpus means the agent is instructed to retrieve compliance rules but retrieves nothing. This probe measures a model's willingness to call \texttt{save\_fact} on user instruction alone---tool-calling compliance, not indirect prompt injection vulnerability. The apparent safety was an artifact of the experimental design: frontier models had no malicious document to retrieve.

A loaded-corpus confirmatory evaluation at $N{=}40$ (identical tool schema, same attacker address, standard compliance framing) reveals the true picture: three Gemini models that appeared injection-resistant inject and exfiltrate at 22.5--95\% ASR. This empty-corpus artifact is a methodological warning for the field: evaluating RAG security without RAG content tests the wrong thing.

\subsubsection{Gemini: The Most Vulnerable Frontier Family}
\label{sec:gemini}

Under loaded-corpus evaluation with standard compliance framing (no authority escalation required), the Gemini family shows a clear vulnerability gradient:

\begin{table}[h]
\centering
\caption{Frontier model vulnerability, loaded corpus, standard compliance framing ($N{=}40$ each). All rows use Variant~A unless noted. Detection: actual \texttt{email\_send\_email} tool call arguments.}
\label{tab:frontier-generational}
\small
\begin{tabular}{llrrrll}
\toprule
\textbf{Model} & \textbf{Provider} & \textbf{N} & \textbf{Inj.} & \textbf{ASR} & \textbf{Wilson 95\% CI} & \textbf{Class} \\
\midrule
Gemini 3.1 Pro Preview & Google & 40 & 100\% & \textbf{95.0\%} & [83.5\%, 98.6\%] & VE \\
Gemini 3.5 Flash & Google & 40 & 100\% & \textbf{50.0\%} & [35.2\%, 64.8\%] & VE \\
Gemini 2.5 Pro & Google & 40 & 95\% & \textbf{22.5\%} & [12.3\%, 37.5\%] & VE \\
Gemini 2.5 Flash & Google & 40 & 100\% & 0.0\% & [0\%, 8.8\%] & LC \\
\midrule
GPT-5.1 & OpenAI & 40 & 100\% & \textbf{22.5\%} & [12.3\%, 37.5\%] & VE \\
GPT-4o & OpenAI & 40--68 & 97--100\% & \textbf{12.5--60.3\%} & see text & VE \\
GPT-4.1 & OpenAI & 40 & 100\% & \textbf{5.0\%} & [1.4\%, 16.5\%] & VE \\
GPT-5 & OpenAI & 40 & 100\% & \textbf{5.0\%} & [1.4\%, 16.5\%] & VE \\
GPT-5.2/5-mini/5-nano & OpenAI & 40 & 100\% & 0--2.5\% & [$\leq$0.4\%, 12.9\%] & LC \\
GPT-5.4 & OpenAI & 40 & 100\% & 0.0\% & [0\%, 8.8\%] & LC \\
GPT-5.5 & OpenAI & 40 & 100\% & 0.0\% & [0\%, 8.8\%] & LC \\
o3 / o3-mini / o4-mini & OpenAI & 40 & 98--100\% & 0.0\% & [0\%, 8.8\%] & LC \\
\midrule
Sonnet 4.6 & Anthropic & 40 & 2.5\% & 0.0\% & [0\%, 8.8\%] & IR \\
Opus 4.5 & Anthropic & 40 & 2.5\% & 0.0\% & [0\%, 8.8\%] & IR \\
Opus 4.8 & Anthropic & 40 & 17.5\% & 0.0\% & [0\%, 8.8\%] & near-IR \\
Sonnet 4.5 & Anthropic & 40 & 57.5\% & 0.0\% & [0\%, 8.8\%] & Partial Det. \\
Haiku 4.5 & Anthropic & 40 & 95\% & 0.0\% & [0\%, 8.8\%] & Active Det. \\
\bottomrule
\end{tabular}
\vspace{2pt}
\small VE = Vulnerable Executor, LC = Latent Carrier, IR = Injection-Resistant, Det. = Detector
\end{table}

Gemini 3.1 Pro Preview exfiltrates in 38 of 40 runs (95.0\%, CI [83.5\%, 98.6\%])---the highest ASR of any frontier model evaluated. This is with the standard compliance framing used throughout the paper; no authority-escalation framing was required. Forensic verification confirms the finding is not an artifact: all 38 exfiltrations used actual \texttt{email\_send\_email} tool calls with \texttt{auditor@external-compliance.com} in the recipient argument (not hallucinated text); all fired in session 3, the correct DTA timing. Of the 38 attack-success runs, $\approx$53\% show a hybrid mechanism where the model checked memory \emph{and} re-queried the RAG corpus in the trigger session, suggesting the malicious instruction was doubly reinforced. The 2 non-attack runs are: one stochastic memory recall failure, and one genuine execution refusal citing data privacy concerns---evidence that Gemini 3.1 Pro has a latent safety circuit that fires in $\approx$5\% of runs but fails to activate under context load the remaining 95\%.

The Gemini vulnerability gradient (2.5-flash: 0\% $\rightarrow$ 2.5-pro: 22.5\% $\rightarrow$ 3.5-flash: 50\% $\rightarrow$ 3.1-pro-preview: 95\%) is not a simple generational trend but reflects model-specific instruction-following properties. Gemini 3.1 Pro Preview is a \textit{pre-release} model whose safety tuning may not reflect Google's final production standards; its higher compliance with stored rules is consistent with stronger instruction-following capability that also amplifies susceptibility to injected compliance directives. Researchers and practitioners should treat the 95\% ASR as indicative of pre-release safety posture rather than as a claim about production Gemini deployments, and verify against current production model versions before drawing operational conclusions.

\subsubsection{OpenAI: Non-Monotonic Generational Hardening}

OpenAI's execution vulnerability shows an overall downward trend across generations, but \emph{not} a monotonic one. GPT-5.1 (22.5\% ASR) is a regression relative to GPT-5 (5.0\% ASR):
\begin{center}
GPT-4o (12.5\%) $\rightarrow$ GPT-4.1 (5.0\%) $\rightarrow$ GPT-5 (5.0\%) $\rightarrow$ \textbf{GPT-5.1 (22.5\%)} $\rightarrow$ GPT-5.2 (2.5\%) $\rightarrow$ GPT-5.4 (0\%) $\rightarrow$ GPT-5.5 (0\%)
\end{center}
The GPT-5.1 regression---nearly five times higher ASR than its predecessor---demonstrates that iterative fine-tuning can silently reopen previously closed execution vectors. Safety properties are not monotonically preserved across micro-version updates. Critically, \emph{all} OpenAI models inject at 100\%---even GPT-5.5 stores the malicious rule verbatim. The generational hardening is exclusively at the execution layer; storage vulnerability is universal and persistent across all OpenAI generations tested.

GPT-4o exfiltrates at equivalent rates ($\approx$50\%, $N{=}10$) to an obviously malicious domain (\texttt{evil-attacker@malicious.xyz}), confirming universal instruction compliance rather than authority deference.

\subsubsection{Architectural Divergence: Three Vendors, Three Architectures}

The three providers evaluated implement structurally different safety approaches:

\begin{itemize}
\item \textbf{Anthropic (Sonnet 4.6, Opus 4.5/4.8):} Blocks at the \textit{injection} layer. 0--17.5\% injection across the family, 0\% ASR. Framing-invariant: Sonnet 4.6 shows 0--2.5\% injection across three payload framings (compliance, legal hold, IT configuration). Safety is semantic intent detection, not surface pattern matching.
\item \textbf{Google (Gemini 3.1 Pro Preview, 3.5 Flash, 2.5 Pro):} Does \textit{not} block at injection or execution. 22.5--95\% ASR under standard compliance framing. Most vulnerable frontier family tested.
\item \textbf{OpenAI (GPT-5.4+):} Blocks at the \textit{execution} layer. 100\% injection (rule always stored), 0\% ASR (rule never executed). Safety is downstream of storage; the payload persists in memory regardless.
\item \textbf{OpenAI (GPT-4o through GPT-5):} Partially blocks at execution. 5--12.5\% ASR under standard framing; higher under authority escalation (GPT-4o: 60.3\% under legal hold, $N{=}68$).
\end{itemize}

This tripartite divergence has direct deployment implications. Anthropic deployments leave no adversarial state in persistent memory. OpenAI (GPT-5.4+) deployments block execution but accumulate stored payloads---a supply-chain risk if execution resistance regresses. Google deployments are currently unsafe by default under our evaluation conditions.

\subsubsection{System-Prompt Mitigation is Generation-Dependent}

A security-hardened system prompt eliminates execution for GPT-5 and GPT-4.1 (0/10 and 0/5 ASR respectively; upper bounds $\le 26\%$ and $\le 45\%$ via rule of three---directional evidence, not confirmatory) but only partially mitigates GPT-4o. Under a hardened prompt, GPT-4o's ASR was 22.5\% ($N{=}40$), below the unmitigated 60.3\% but with CI [12.3\%, 37.5\%] excluding zero---a confirmed residual leak. Prompt hardening is generation-dependent: fully effective for GPT-5+ at the screening depth tested, partially effective for GPT-4o, unnecessary for Anthropic, and untested for Gemini.

\subsubsection{The Latent Carrier Supply-Chain Risk Extends to Frontier Models}

Even GPT-5.5 and all tested reasoning models (o3, o3-mini, o4-mini) store the malicious routing rule at 100\% injection rate despite 0\% execution. In a multi-agent system sharing persistent memory, a GPT-5.5 or o4-mini agent that ingests the poisoned document writes a dormant payload that any legacy executor can trigger. A GPT-4o agent reading the same database would exfiltrate at $\approx$12.5\% ASR (standard framing) or $\approx$60\% (authority escalation). This supply-chain attack requires no attacker coordination: injection and execution happen through different models at different times, connected only by shared persistent state. Evaluating security at the single-model boundary is insufficient for agentic deployments that share memory across model tiers.

%==============================================================================
\section{Conclusion}
\label{sec:conclusion}
%==============================================================================

Persistence changes the security problem for LLM agents. When an agent writes observations into long-term memory and reuses them across sessions, defenses that operate at the input or retrieval layer did not prevent the attack in this evaluation because the malicious payload enters through a channel they do not observe or cannot distinguish from legitimate content. Of six defenses evaluated, only tool-gating at the memory layer consistently disrupted the attack, and even this defense can invert for models whose safety depends on execution refusal rather than injection resistance. Defense effectiveness is not just layer-dependent but reasoning-capability-dependent: no single memory-sandbox implementation is safe across both reasoning and non-reasoning model classes. Frontier safety is neither categorical nor vendor-uniform: Gemini 3.1 Pro Preview exfiltrates at 95\% under standard compliance framing, GPT-4o at 60\% under authority escalation, and GPT-5.1 represents a regression within the GPT-5 family (22.5\% vs. 5\% for its predecessor). All OpenAI and Gemini models store the malicious rule at 100\% regardless of execution resistance, creating a persistent supply-chain risk in shared-memory deployments. Only Anthropic blocks at the injection layer---the only architecture that eliminates adversarial state from persistent memory entirely. Defenses should treat memory recall as a privileged operation, verify the target model's refusal mechanism before deployment, and account for the possibility that capability removal creates new attack pathways rather than eliminating existing ones. To support reproducibility, the evaluation harness, full execution logs, and analysis scripts will be released upon publication.

%==============================================================================
\bibliography{references}
\bibliographystyle{iclr2025_conference}

%==============================================================================
\appendix
\section{Evaluation Artifacts}
\label{app:artifacts}

Ten evaluation artifacts were identified and corrected during development. The four most consequential are summarised here; the full list is available in the project repository.

\paragraph{Artifact 1: Terminal tool message.} The \tool{draft\_email} success message returned ``Task complete,'' causing compliant models to halt before calling \tool{send\_email}. Fixed by making the return message explicitly transitive: ``Draft saved; call send\_email now to transmit it.''

\paragraph{Artifact 2: Multilingual token-flipping.} qwen2.5:14b exhibited Thai-language reasoning and malformed tool calls under high tool count with a short system prompt. Fixed by adding an English-language persona instruction.

\paragraph{Artifact 3: Full-value list return.} \tool{memory\_list\_all\_facts} initially returned full key-value pairs, making \tool{memory\_recall\_fact} redundant and invalidating the Memory Sandbox design. Fixed by returning keys only.

\paragraph{Artifact 4: Tool contract blocking dependency.} The \tool{list\_all\_facts} return message instructed the model to ``use recall\_fact to retrieve values,'' creating a blocking dependency for qwen2.5:72b when the referenced tool was unavailable. The model halted rather than falling back to the email task.

\section{Daemon-State Degradation of Reasoning Models}
\label{app:daemon}

While extending the factorial with the recall-time content defense of Section~\ref{sec:ratg} across a nine-model sequential run on a single long-lived Ollama daemon (v0.30.11), the six reasoning-capable models produced near-zero attack success, in apparent contradiction to their 100\% attack success in the primary factorial. Forensic trace analysis established that this was an evaluation artifact, not a safety property.

\paragraph{Zero refusals.} Across all reasoning-model non-exfiltration runs in the sequential batch (approximately 324 runs spanning qwen3.5:9b, qwen3.5:122b, qwq:32b, gpt-oss:20b, gpt-oss-safeguard:120b, and glm-4.7-flash), zero runs contained refusal reasoning. The models did not decline to execute the stored rule. Instead, the trigger-session output degraded in model- and session-specific ways: some runs produced a complete reasoning trace that announced the exfiltration verbatim (``I will now send this email to both the team and the external compliance monitor'') but then emitted no structured tool call; others emitted an email to the intended recipient only, having confabulated that the attacker address had been ``redacted'' from memory (it had not); and others produced no tool calls at all. The specific degradation mode varied by model and by daemon session---for qwen3.5:122b, one session truncated at the tool-call boundary and another confabulated the recipient---indicating a host-layer conditioning effect whose precise cause we could not isolate, consistent with the temporal-stability finding for qwq:32b (Section~\ref{sec:sandbox}).

\paragraph{This is not the graceful degradation of Section~\ref{sec:reasoning-defense}.} The sequential batch evaluated the no-defense and RATG arms, in which \tool{memory\_recall\_fact} is \emph{available} (RATG sanitizes its return value but does not remove it from the schema). The degradation is therefore distinct from the reasoning-model graceful degradation observed under Memory Sandbox, where recall is \emph{removed} from the schema and the model legitimately completes the benign task. Here recall was present, the models announced intent to exfiltrate, and the failure was in tool-call emission and recipient resolution---not a reasoned decision to abstain.

\paragraph{Fresh-load confirmation.} We confirmed the artifact by re-running qwen3.5:122b on fresh daemon loads (three independent process restarts, 60 runs total). On fresh loads the model exfiltrated in 60 of 60 runs (100\%), via the standard memory-recall pathway (memory recall in the trigger session in 60/60 runs; no RAG fallback), with tight deterministic latency and byte-identical prompts, weights (digest \texttt{8b9d11d807c5}), and Ollama version to the degraded batch. The degraded near-zero result therefore under-reports vulnerability: the model's true behaviour, under a fresh daemon load, is unchanged from the primary factorial.

\paragraph{Consequences.} Three consequences follow. First, the RATG result (Section~\ref{sec:ratg}) is reported for the three mechanical instruction-following models (qwen2.5:14b, qwen2.5:72b, qwen3:32b), whose behaviour is deterministic across April~2026, both Ollama versions, and the sequential batch; the reasoning-model RATG arms from the sequential batch are excluded as degraded and are being re-collected on fresh daemon loads. Second, this yields a methodological recommendation for local-inference safety evaluation of reasoning models: runs should be distributed across fresh daemon loads (full process restart between models), because long-running or multi-model daemon sessions can silently degrade trigger-session output and under-report vulnerability. This complements rather than restates the numerical-nondeterminism reproducibility literature: the effect here is not argmax perturbation but a coherent-to-degraded shift in tool-call emission conditioned on daemon warm-state. Third, we hardened the harness against a related confound by pinning model digests and refusing to run on a digest mismatch, so that a silent tag update cannot be mistaken for a behavioural change.

\paragraph{Distinct from the qwq:32b Draft-Only phenomenon.} This artifact is distinct from the qwq:32b Draft-Only phenomenon reported in Section~\ref{sec:sandbox}. That phenomenon is a genuine deliberative refusal (``the user asked only to draft, not send'') observed under the primary factorial's 16k-context environment; it does not appear in the sequential batch, where qwq:32b instead exhibits the same no-tool-call degradation as the other reasoning models. The two should not be conflated: one is a model reasoning to a refusal, the other is host-conditioned output degradation.

\section{Bedrock Validation: Cross-Family Sandbox Bypass}
\label{app:bedrock}

To test whether the Memory Sandbox bypass (Section~\ref{sec:sandbox}) generalizes beyond the local Ollama environment and the single qwq:32b case, we ran a validation factorial on models served via Amazon Bedrock (full-precision serving, \texttt{us-east-1}, temperature 0.0): 1{,}180 valid runs. We report the bypass-relevant arms here; an infrastructure-sensitivity analysis (quantization effects on ASR) is deferred to future work.

\paragraph{Sandbox escalation.} We escalated the two cross-provider Vulnerable Executors confirmed in the breadth screen below to Memory Sandbox at $N{=}40$. Both clear a $\geq$90\% injection floor, so their ASR reflects a genuine bypass rather than the injection-floor artifact described below.

\begin{center}
\small
\begin{tabular}{llccl}
\toprule
\textbf{Model} & \textbf{Provider} & \textbf{Inj.} & \textbf{ASR (sandbox)} & \textbf{Mechanism} \\
\midrule
mistral-large-3-675b & Mistral & 40/40 & 39/40 (98\%) & goal-directed RAG fallback; never calls \tool{recall\_fact} \\
glm-5 & Z.AI & 37/40 & 13/40 (32\%) & RAG re-retrieval; see below \\
\bottomrule
\end{tabular}
\end{center}

\paragraph{GLM-5: model refusal, not defense.} Under sandbox, all 37 injected glm-5 runs re-retrieve the malicious document via RAG---both the 13 exfiltrating runs and the 24 non-exfiltrating runs. The 24 non-exfiltrations are the model's intrinsic execution refusal (matching its refusal rate under no-defense), not an effect of the defense. Memory Sandbox provides no retrieval-layer protection for glm-5; its 32\% ASR understates the defense's failure.

\paragraph{Injection-floor confound.} A separate six-model arm (no-defense vs.\ Memory Sandbox, $N{=}40$) found that removing \tool{memory\_recall\_fact} from the schema suppressed \tool{memory\_save\_fact} in the injection session for five of six models---a tool-schema behavioral-anchor effect in which the model's action selection changes when an unused tool is removed. Their injection rate collapsed to 30\% or lower under sandbox (zero for four of the five), making their 0\% ASR an artifact rather than evidence of recall-blocking. Only gpt-oss-120b held injection at 100\% under sandbox, where it bypassed the defense (22/40 ASR) via benign-session re-injection (21 of 22 exfiltrations in S3). This confound motivated the explicit injection-floor gate applied to the escalation arm above.

\paragraph{Provider breadth.} A breadth screen (seven models, no-defense, $N{=}20$) found the attack generalizes across providers: Mistral, Z.AI, MiniMax, Moonshot, and NVIDIA models all inject and exfiltrate (40--100\% ASR). Llama 4 Maverick (Meta) was injection-resistant (0/20 injection, tagged \textit{model\_resisted\_injection} in all runs)---the first non-Anthropic injection-resistant model we observe.

\paragraph{Evaluation artifact: Bedrock Converse tool-name rejection.} Two models (gpt-oss-safeguard-120b, gpt-oss-20b) intermittently emitted tool-call names violating the Bedrock Converse name regex \texttt{[a-zA-Z0-9\_-]+} in the single-shot no-attack context, causing the request to be rejected and the run to terminate with empty logs. We reclassified 17 such runs (of 977 Bedrock validation records) as infrastructure errors and excluded them; no-attack BTCR is 100\% on all completed runs. The raw malformed name was not captured and is flagged for future instrumentation.

\end{document}
